% 附录 D IC Validator PXL Debugger
\biappendix{IC Validator PXL Debugger}{IC Validator PXL Debugger}
\label{app:d}

IC Validator PXL Debugger 协助你构造 runset 或 runset 的一部分，使你能够查看正在编写的每条设计规则或器件配置函数中的中间阶段。

PXL Debugger 包括以下各节：
\begin{itemize}
  \item 运行 PXL Debugger（Running the PXL Debugger）
  \item PXL Debugger 命令（PXL Debugger Commands）
\end{itemize}

% ============================================================
\section{运行 PXL Debugger}
\subsection*{Running the PXL Debugger}

要运行 PXL Debugger，在 IC Validator 命令中使用 \cmd{-debug} 命令行选项：

\begin{lstlisting}
% icv -debug runset.rs [-debug_source file]
\end{lstlisting}

\cmd{-debug\_source} 命令行选项允许以批处理模式运行 PXL Debugger。指定先前调试运行中用 \cmd{record} 命令写出的文件。PXL Debugger 的输出显示在屏幕上。

当 runset 执行开始时：
\begin{itemize}
  \item 屏幕消息指示 PXL Debugger 已启动；
  \item 运行在第一个断点处暂停。显示可执行行之前五行与之后四行；
  \item 显示 PXL Debugger 提示符（\cmd{ICVD>>}）。
\end{itemize}

例如：

\begin{lstlisting}
Starting ICV Debugger
       13 |          t = t + dtoi((x * y**2) / z);
       14 |      }
       15 |      t = dtoi(t / x);
       16 |      note (" f2(3) = "+ t);
       17 | }
  --> 18 | xx:integer;
       19 | f2(3);
       20 | note("end. ");
ICVD>>
\end{lstlisting}

在提示符处输入希望 PXL Debugger 执行的操作。见下文「PXL Debugger 命令」。

\begin{noteBox}
PXL Debugger 设计为顺序执行 PXL 代码。若启用了分布式处理选项，将生成警告并忽略这些选项。

PXL Debugger 在 runset 与用户函数中单步执行代码，但不在远程（remote）函数中单步执行。
\end{noteBox}

% ============================================================
\section{PXL Debugger 命令}
\subsection*{PXL Debugger Commands}

IC Validator PXL Debugger 的命令见表~\ref{tab:app-d-cmds}。$^{a}$

\begin{longtable}{@{}>{\ttfamily\raggedright\arraybackslash}p{0.32\textwidth} p{0.62\textwidth}@{}}
\caption{PXL Debugger 命令}\\
\label{tab:app-d-cmds}\\
\toprule
\textbf{\rmfamily 命令} & \textbf{定义} \\
\midrule
\endfirsthead
\toprule
\textbf{\rmfamily 命令} & \textbf{定义} \\
\midrule
\endhead
\bottomrule
\endfoot

h[elp] & 打印 PXL Debugger 命令列表。 \\
\midrule
s[tep] & 步入函数调用。见 \cmd{step} 命令说明。 \\
\midrule
n[ext] & 一次性执行完函数。见 \cmd{next} 命令说明。 \\
\midrule
q[uit] & 退出 PXL Debugger 与 IC Validator 运行。见 \cmd{quit} 命令说明。 \\
\midrule
c[ont] & 继续直至下一断点或运行结束。见 \cmd{cont} 命令说明。 \\
\midrule
fin[ish] & 执行直至当前函数返回。见 \cmd{finish} 命令说明。 \\
\midrule
b[reak]\newline
b[reak] [\textit{file}:]\textit{line}\newline
b[reak] \textit{function} &
设置断点。默认为当前文件:行。见 \cmd{break} 命令说明。 \\
\midrule
info b[reak] [\textit{number}] &
显示断点状态。默认为全部断点。见 \cmd{info break} 命令说明。 \\
\midrule
dis[able] [\textit{number}] &
禁用断点。默认为禁用全部断点。见 \cmd{disable} 命令说明。 \\
\midrule
en[able] [\textit{number}] &
启用断点。默认为启用全部断点。见 \cmd{enable} 命令说明。 \\
\midrule
d[elete] \textit{number} &
删除指定断点。见 \cmd{delete} 命令说明。 \\
\midrule
l[ist]\newline
l[ist] [\textit{file}:]\textit{line}\newline
l[ist] \textit{function} &
列出源码。默认为当前文件与行。见 \cmd{list} 命令说明。 \\
\midrule
where & 显示函数调用栈。见 \cmd{where} 命令说明。 \\
\midrule
pwd &
打印 IC Validator PXL Debugger 的当前工作目录。例如：
\begin{lstlisting}[basicstyle=\ttfamily\footnotesize]
ICVD>> pwd
Working directory /remote/data/ICV
\end{lstlisting} \\
\midrule
p[rint] \textit{expression} &
打印表达式的值。见 \cmd{print} 命令说明。 \\
\midrule
record [\textit{file}] &
保存当前调试会话的命令。默认为 \cmd{ICV\textit{pid}.pdi}。见 \cmd{record} 命令说明。 \\
\midrule
source \textit{file} &
从文件执行命令。见 \cmd{source} 命令说明。 \\
\end{longtable}

{\footnotesize
\begin{itemize}[leftmargin=1.5em,nosep]
  \item[$^{a}$] 命令可按方括号所示缩写。
\end{itemize}
}

\subsection{break}
\subsubsection*{break}

\cmd{break} 命令设置断点。有四种方式：

\begin{enumerate}
  \item \cmd{break}（无选项）：在当前文件的当前行设置断点。
\begin{lstlisting}
        4 |
        5 | f2 : function (
        6 |     x : integer
        7 |     ) returning ans : integer
        8 | {
  --> 9 |        ans = x;
       10 |      note("0. ans = " + ans);
       11 | }
       12 |
       13 | xx: integer = 0;
ICVD>> b
Breakpoint #1, f2() at test.rs:9
ICVD>>
\end{lstlisting}

  \item \cmd{break} \textit{line}：在当前文件的指定行号设置断点。
\begin{lstlisting}
ICVD>> b 13
Breakpoint #1,      f2() at test.rs:13
ICVD>>
\end{lstlisting}

  \item \cmd{break} \textit{file}:\textit{line}：在指定文件的指定行号设置断点。
\begin{lstlisting}
ICVD>> b test.rs:15
Breakpoint #2, _MAIN_() at               test.rs:15
ICVD>>
\end{lstlisting}

  \item \cmd{break} \textit{function}：在指定函数内第一条可执行行设置断点。
\begin{lstlisting}
       8 | {
       9 |      ans = x;
      10 |      note("0. ans = " + ans);
      11 | }
      12 |
  --> 13 | xx: integer = 0;
      14 | xx = f2(3);
      15 | note("1. f2(3)= " + xx);
ICVD>> b f2
Breakpoint #1,      f2() at test.rs:9
ICVD>>
\end{lstlisting}

此方法可在重载函数上设置断点。下例中 \cmd{f2} 可返回 integer 或 void。指定 2 选择选项 2（返回 integer）：
\begin{lstlisting}
ICVD>> b f2
[0] cancel
[1] all
[2] f2() returning integer, at test.rs:9
[3] f2() returning void, at test.rs:16
> 2
\end{lstlisting}

指定 1 选择选项 1（返回 integer 与 void）：
\begin{lstlisting}
Breakpoint #1, f2() at test.rs:9
ICVD>> b f2
[0] cancel
[1] all
[2] f2() returning integer, at test.rs:9
[3] f2() returning void, at test.rs:16
> 1
Breakpoint #2, f2() at test.rs:9
Breakpoint #3, f2() at test.rs:16
ICVD>>
\end{lstlisting}
\end{enumerate}

\subsection{cont}
\subsubsection*{cont}

\cmd{cont} 命令执行 runset 直至完成，或在指定断点处停止。

下例中 runset 完成调试运行：
\begin{lstlisting}
       1 | #include <math.rh>
       2 | #include <diagnostics.rh>
       3 |
  --> 4 | f: list of integer = {1, 1};
       5 |
       6 | for (i in 2 thru 19) {
       7 |        f.push_back(f[i-2] + f[i-1]);
       8 | }
ICVD>> c

Completed ICV Debugger Run
ICVD>>
\end{lstlisting}

下例中执行停在循环内的断点：
\begin{lstlisting}
ICVD>> b 13
Breakpoint #1,       f2() at test.rs:13
ICVD>> c
        8 | {
        9 |          t : integer = 0;
       10 |          y : integer = 10;
       11 |          z : integer = 100;
       12 |          for ( i=1 to x) {
B --> 13 |               t = t + dtoi((x * y**2) / z);
       14 |          }
       15 |          t = dtoi(t / x);
       16 |          note (" f2(3) = "+ t);
       17 | }
Breakpoint #1,       f2() at test.rs:13
ICVD>> c
...
Breakpoint #1,       f2() at test.rs:13
ICVD>>
\end{lstlisting}

\subsection{delete}
\subsubsection*{delete}

\cmd{delete} 命令删除指定断点。格式为：

\begin{lstlisting}
delete number
\end{lstlisting}

\begin{lstlisting}
ICVD>> info b
Breakpoint #1,       _MAIN_() at    test.rs:7       Enabled
Breakpoint #2,       _MAIN_() at    test.rs:16      Enabled
ICVD>> d 1
ICVD>> info b
Breakpoint #2,       _MAIN_() at    test.rs:16      Enabled
ICVD>>
\end{lstlisting}

\subsection{disable}
\subsubsection*{disable}

\cmd{disable} 命令禁用指定断点。有两种方式：

\begin{enumerate}
  \item \cmd{disable}（无断点号）：禁用全部断点。
\begin{lstlisting}
ICVD>> info b
Breakpoint #1,      _MAIN_() at          test.rs:7        Enabled
Breakpoint #2,      _MAIN_() at          test.rs:18       Enabled
ICVD>> dis
ICVD>> info b
Breakpoint #1,      _MAIN_() at          test.rs:7        Disabled
Breakpoint #2,      _MAIN_() at          test.rs:18       Disabled
\end{lstlisting}

  \item \cmd{disable} \textit{number}：仅禁用该断点。
\begin{lstlisting}
ICVD>> dis 1
ICVD>> info b
Breakpoint #1,      _MAIN_() at          test.rs:13       Disabled
Breakpoint #2,      _MAIN_() at          test.rs:16       Enabled
ICVD>>
\end{lstlisting}
\end{enumerate}

\subsection{enable}
\subsubsection*{enable}

\cmd{enable} 命令启用先前禁用的断点。有两种方式：

\begin{enumerate}
  \item \cmd{enable}（无断点号）：启用全部断点。
\begin{lstlisting}
ICVD>> info b
Breakpoint #1,      _MAIN_() at          test.rs:7        Disabled
Breakpoint #2,      _MAIN_() at          test.rs:18       Disabled
ICVD>> en
ICVD>> info b
Breakpoint #1,      _MAIN_() at          test.rs:7        Enabled
Breakpoint #2,      _MAIN_() at          test.rs:18       Enabled
ICVD>>
\end{lstlisting}

  \item \cmd{enable} \textit{number}：仅启用该断点。
\begin{lstlisting}
ICVD>> info b
Breakpoint #1,      _MAIN_() at          test.rs:7        Disabled
ICVD>> en 1
ICVD>> info b
Breakpoint #1,      _MAIN_() at          test.rs:7        Enabled
ICVD>>
\end{lstlisting}
\end{enumerate}

\subsection{finish}
\subsubsection*{finish}

\cmd{finish} 命令执行到所在函数结束。但若遇到断点仍会停止。

下例中运行继续直至当前函数完成：
\begin{lstlisting}
       8 | {
       9 |      t : integer = 0;
      10 |      y : integer = 10;
      11 |      z : integer = 100;
      12 |      for ( i=1 to x) {
  --> 13 |          t = t + dtoi((x * y**2) / z);
      14 |      }
      15 |      t = dtoi(t / x);
      16 |      note (" f2(3) = "+ t);
      17 | }
ICVD>> fin
test.rs:16 f2(3) = 3
      15 |      t = dtoi(t / x);
      16 |      note (" f2(3) = "+ t);
      17 | }
      18 | xx:integer;
      19 | f2(3);
  --> 20 | note("end. ");
\end{lstlisting}

下例中运行继续直至遇到断点：
\begin{lstlisting}
ICVD>> fin
test.rs:11 fibonacci(16) = 1597
...
B --> 11 |       note("fibonacci(" + i + ") = " + f[i]);
...
Breakpoint #2, _MAIN_() at       test.rs:11
ICVD>> fin
test.rs:11 fibonacci(17) = 2584
...
Breakpoint #2, _MAIN_() at       test.rs:11
ICVD>>
\end{lstlisting}

\subsection{info break}
\subsubsection*{info break}

\cmd{info break} 命令列出本 PXL Debugger 会话中设置的断点信息。有两种方式：

\begin{enumerate}
  \item \cmd{info break}（无断点号）：列出全部断点。
\begin{lstlisting}
ICVD>> info b
Breakpoint #1,       _MAIN_() at          test.rs:7         Enabled
Breakpoint #2,       _MAIN_() at          test.rs:16        Enabled
\end{lstlisting}

  \item \cmd{info break} \textit{number}：仅列出该断点信息。
\begin{lstlisting}
ICVD>> info b 2
Breakpoint #2, _MAIN_() at                test.rs:16        Enabled
ICVD>>
\end{lstlisting}
\end{enumerate}

\subsection{list}
\subsubsection*{list}

\cmd{list} 命令列出源码行。有四种方式：

\begin{enumerate}
  \item \cmd{list}（无选项）：从上次打印的行起列出 10 行。
\begin{lstlisting}
        1 | #include <math.rh>
        2 | #include <diagnostics.rh>
        3 |
  --> 4 | f: list of integer = {1, 1};
        5 |
        6 | for (i in 2 thru 19) {
        7 |        f.push_back(f[i-2] + f[i-1]);
        8 | }
ICVD>> list
        9 |
       10 | /* comment 1
...
       17 | for (i in 0 thru 19) {
       18 |      note("fibonacci(" + i + ") = " + f[i]);
ICVD>>
\end{lstlisting}

  \item \cmd{list} \textit{line}：从当前文件打印指定行号之前五行与之后四行。

  \item \cmd{list} \textit{file}:\textit{line}：从指定文件打印指定行号之前五行与之后四行。
\begin{lstlisting}
ICVD>>   list test.rs:4
       1 | #include <math.rh>
       2 | #include <diagnostics.rh>
       3 |
  --> 4 | f: list of integer = {1, 1};
       5 |
       6 | for (i in 2 thru 19) {
       7 |        f.push_back(f[i-2] + f[i-1]);
       8 | }
ICVD>>
\end{lstlisting}

  \item \cmd{list} \textit{function}：打印指定函数内第一条可执行行之前五行与之后四行。
\begin{lstlisting}
ICVD>>    list f2
         4 |
         5 | f2 : function (
         6 |      x : integer
         7 |      ) returning void
         8 | {
         9 |       t : integer = 0;
...
ICVD>>
\end{lstlisting}
\end{enumerate}

\subsection{next}
\subsubsection*{next}

\cmd{next} 命令将 PXL Debugger 提示符移到 runset 中下一可执行行。对函数调用执行 \cmd{next} 时，会完成该函数执行，再移到 runset 中下一可执行行。

在 runset 可执行行上给出 \cmd{next}：
\begin{lstlisting}
       13 |          t = t + dtoi((x * y**2) / z);
...
  --> 18 | xx:integer;
       19 | f2(3);
       20 | note("end. ");
ICVD>> n
...
  --> 19 | f2(3);
       20 | note("end. ");
ICVD>>
\end{lstlisting}

在函数调用上给出 \cmd{next}：
\begin{lstlisting}
  --> 13 | xx = f2(3);
       14 | note("1. f2(3)= " + xx);
ICVD>> n
...
  --> 14 | note("1. f2(3)= " + xx);
ICVD>>
\end{lstlisting}

\subsection{print}
\subsubsection*{print}

\cmd{print} 命令打印变量与简单表达式的值。格式为：

\begin{lstlisting}
print expression
\end{lstlisting}

打印变量值：
\begin{lstlisting}
ICVD>> p ans
ans = 3
\end{lstlisting}

打印循环变量：
\begin{lstlisting}
ICVD>> p i
i = 3
\end{lstlisting}

打印表达式：
\begin{lstlisting}
ICVD>> p xx
xx = 0
ICVD>> p (xx+7)*20
(xx+7)*20 = 140
ICVD>>
\end{lstlisting}

\subsection{quit}
\subsubsection*{quit}

\cmd{quit} 命令结束调试会话与 IC Validator 运行。

\begin{lstlisting}
ICVD>> q
IC Validator Run: Time= 0:00:04 User=0.06 System=0.01
Overall icv_engine time = 0:00:04
IC Validator is done.
\end{lstlisting}

\subsection{record}
\subsubsection*{record}

\cmd{record} 命令将当前调试会话的命令保存到 ASCII 文件。有两种方式：

\begin{enumerate}
  \item \cmd{record}（未指定文件名）：命令保存到 \cmd{ICV\textit{pid}.pdi}，其中 \textit{pid} 为 UNIX 进程 ID。
\begin{lstlisting}
ICVD>> record
Saving Commands to 'ICV22347.pdi'
ICVD>>
\end{lstlisting}

  \item \cmd{record} \textit{file}：可指定带或不带路径的文件名。
\begin{lstlisting}
ICVD>> record myfile.cmd
Saving Commands to 'myfile.cmd'
ICVD>>
\end{lstlisting}
\end{enumerate}

ASCII 文件 verbatim 存储除 \cmd{break} 与 \cmd{list} 外的全部命令；\cmd{break} 与 \cmd{list} 以解析后的文件路径存储。例如会话后记录文件含：

\begin{lstlisting}
b f2
list f2
b b.rs:6
q
\end{lstlisting}

\subsection{source}
\subsubsection*{source}

\cmd{source} 命令执行指定文件中的 PXL Debugger 命令。命令完成后，控制返回用户，除非源文件中的 \cmd{quit} 使调试器结束。格式为：

\begin{lstlisting}
source file
\end{lstlisting}

例如：
\begin{lstlisting}
ICVD>>   source ICVDebug9398.pdi
ICVD>>   n
...
ICVD>> b 7
Breakpoint #1, _MAIN_() at       test.rs:7
ICVD>> b 16
Breakpoint #2, _MAIN_() at       test.rs:16
ICVD>> c
...
Breakpoint #1, _MAIN_() at       test.rs:7
End of commands from commandfile
ICVD>>
\end{lstlisting}

\subsection{step}
\subsubsection*{step}

\cmd{step} 命令允许步入函数调用，执行函数内每一行。光标移到函数内第一条可执行行。

下例中步入 \cmd{f2}，第一条可执行行为 9：
\begin{lstlisting}
  --> 19 | f2(3);
       20 | note("end. ");
ICVD>> s
        4 |
        5 | f2 : function (
...
  --> 9 |        t : integer = 0;
...
ICVD>>
\end{lstlisting}

\subsection{where}
\subsubsection*{where}

\cmd{where} 命令打印 runset 中某语句的函数调用栈。下例显示已有两次函数调用：

\begin{lstlisting}
ICVD>> where
#0     | sum()     at test.rs:27
#1     | _MAIN_()          at test.rs:32
\end{lstlisting}

下例显示已有三次函数调用：

\begin{lstlisting}
ICVD>> where
#0     | add()     at test.rs:22
#1     | sum()     at test.rs:27
#2     | _MAIN_()          at test.rs:32
ICVD>>
\end{lstlisting}
